% !PRINTED 212
\weberpage{212}
Irgend eine rationale Gleichung zwischen den Wurzeln $x_{\mu,\mu'}$ und $\varkappa^2$ geht, auch wenn beliebige konstante, d.\,h. von $\omega$ unabhängige Zahlenkoeffizienten darin vorkommen, durch (10) in eine Identität über, und man kann daher für $\omega$

\[
\frac{\gamma + \delta\omega}{\alpha + \beta\omega}
\]

substituieren, d.\,h. man kann jede Substitution $T$ auf die rationale Gleichung zwischen den $x_{\mu,\mu'}$ anwenden. Daraus folgt, \gesperrt{daß die ganze Gruppe $\mathfrak{B}$ in $\mathfrak{G}$ enthalten ist (Bd.\,I, § 156) und sogar in der Gruppe $\mathfrak{G}'$, die aus der Gruppe $\mathfrak{G}$ der Teilungsgleichung durch Adjunktion beliebiger Zahlen entsteht.}

Wenn wir den Rationalitätsbereich der rationalen Zahlen durch Adjunktion von $n$ten Einheitswurzeln erweitern, so gehören zu den rationalen Gleichungen auch die \gesperrt{Abelschen} Relationen des vorigen Paragraphen. Auf diese Relationen, etwa auf

\[
\overset{\nu}{\Sigma} e^{\frac{8\nu\nu'\pi i}{n}} x_{\nu,\nu'} = 0
\]

ist aber \gesperrt{keine} der Substitutionen

\[
M = \begin{pmatrix} m, & 0 \\ 0, & 1 \end{pmatrix},
\]

in der $m$ von 1 verschwindet, % [sic] printed \textquotedbl{}verschwindet\textquotedbl{}; expected sense \textquotedbl{}verschieden ist\textquotedbl{} (differs from 1)
anwendbar; denn durch diese Substitution geht, wenn $mm' \equiv 1 \pmod{n}$ ist,

\[
\overset{\nu}{\Sigma} e^{\frac{8\nu\nu'\pi i}{n}} x_{\nu,\nu'} \ \text{in} \ \overset{\nu}{\Sigma} e^{\frac{8\nu m'\nu'\pi i}{n}} x_{\nu,\nu'}
\]

über, was nach (12), § 62 von Null verschieden ist, wenn nicht $m'$ und also auch $m \equiv 1 \pmod{n}$ ist. Daraus folgt der Satz:

2. \gesperrt{Nach Adjunktion der $n$ten Einheitswurzeln (und beliebiger anderer Konstanten) ist $\mathfrak{B}$ die Galoissche Gruppe der Teilungsgleichung. Es wird $\mathfrak{B}$ auch die Monodromiegruppe der Teilungsgleichung genannt.}

3. Wir beweisen jetzt noch, daß in dem ursprünglichen Rationalitätsbereich, also ohne Adjunktion der $n$ten Einheitswurzeln oder anderer irrationaler Zahlen, \gesperrt{die Gruppe der Teilungsgleichung mit der Gruppe $\mathfrak{A}$ identisch ist.}

Es ist in Bd.\,I, § 174 gezeigt, daß die primitiven $n$ten Einheitswurzeln $\varrho$ Wurzeln einer irreduziblen Gleichung

\weq{12}{\Phi(\varrho) = 0}
%CONT>>
